% Separate continuation: the original norm-image calculation.
% GPT-6 Astra Pro; /root/fixed_surface; 2026-11-22 UTC.
% Positivity and irregularity argument developed with /root.

\section{HD-F1-S10: semiregularity of the original norm-image surface}
\label{sec:f1-norm-image-semiregularity}

Consider the actual images $S_j\subset P$ of HD-F1-S09 and HD-F1-T07,
after translating each connected norm torsor to its Prym fourfold.
For every connected etale cyclic triple cover of a smooth
nonhyperelliptic genus-three curve, the following statements hold:
\[
 \begin{gathered}
 T_j\longrightarrow S_j\text{ contracts exactly one complete pencil},\\
 K_{S_j}^2=16,\quad c_2(S_j)=8,\quad q(S_j)=4,\quad p_g(S_j)=5,\\
 [S_j]=\beta^2/6+w_j,\quad 0\ne w_j\in W_K(P),\quad w_j^2=2,\\
 \operatorname{rank}\mu_{S_j}=12,\qquad
 \dim\ker\mu_{S_j}=4,\qquad\operatorname{coker}\mu_{S_j}=0.
 \end{gathered}
 \tag{F1S10.1}
\]
Here $\beta=\Theta_Z|_P$ is the specified polarization of type
$(1,1,3,3)$, $W_K(P)=\bigwedge_K^4 H^1(P,\mathbb Q)$, and
$\mu$ is precisely the conormal map of DEG\_L S-02.
The proof below uses the original norm surface throughout. Paragraphs
1--5 establish the calculation at a very general point of each cover
component. Paragraph 6 proves its extension to every member just stated.

\paragraph{1. The generic Hodge space and the nonzero determinant component.}
Schoen (1988), the differential calculation in Lemma 3.11 and the
argument immediately following it, proves that the Prym period map
for these covers dominates its four-dimensional unitary family.
The displayed multiplication map has source
$H^0(K_C\otimes\epsilon)\otimes H^0(K_C\otimes\epsilon^{-1})$
of dimension four, and is injective at a general base curve by the
base-point-free pencil trick. Here $\epsilon$ is the nontrivial
three-torsion bundle defining the cover. This assertion concerns
this particular Prym component; no assertion about other unitary
components is required.

At a very general point of this domain, $P$ is simple and
\[
 H^4(P,\mathbb Q)\cap H^{2,2}(P)
                     =\mathbb Q\beta^2\oplus W_K(P).
 \tag{F1S10.2}
\]
One can verify the generic assertion directly from the domain.
The complexification of $H^1$ is $V\oplus V^*$ for the standard
four-dimensional representation of $\mathrm{SL}_4$.
The conjugates of the noncentral Hodge circle generate the
unitary derived group, whose complexification is this
$\mathrm{SL}_4$. The invariants in
$\bigwedge^4(V\oplus V^*)$ have dimension three: the two
determinants $\bigwedge^4V$, $\bigwedge^4V^*$, and the scalar
in $\bigwedge^2V\otimes\bigwedge^2V^*$. These are precisely
the three classes in (F1S10.2). Every additional rational
Hodge class lies on its own proper Hodge locus; there are
only countably many rational classes. The same argument
for endomorphisms gives generic endomorphism algebra $K$.
As a field this has no nontrivial idempotent, so $P$ is simple.

We verify that the norm images themselves have a nonzero
component in $W_K$, including the maps from Schoen's cycles.
Choose an Abel map $a_Z:Z\to JZ$ and define
\[
 Q_0=3\operatorname{Id}-(\operatorname{Id}+\sigma+\sigma^2),
 \qquad
 f:Z^4\xrightarrow{\sum a_Z}JZ\xrightarrow{Q_0}P.
 \tag{F1S10.3}
\]
The curve $Q_0a_Z(Z)$ generates $P$, since $Q_0:JZ\to P$
is surjective. Its fourth sum therefore dominates $P$:
successive sums increase dimension until dimension four,
since a dimension stagnation would make the sum invariant
under the curve and hence under the abelian variety it
generates. Thus $f$ is generically finite of positive degree.

Let $U\subset H^4(Z^4,\mathbb Q)$ be Schoen's two-dimensional
primitive-character subspace. Then
\[
                         f^*W_K(P)=U.
 \tag{F1S10.4}
\]
To check this map explicitly, pull back the four one-forms
in either determinant eigenline and expand along the four
factors. Two forms of the same nontrivial character wedge
to zero on one curve, since their product character is
nontrivial whereas $H^2(Z)$ is trivial under the deck group.
Only terms with one form on each factor remain. Their
alternating tensors are exactly Schoen's determinant
eigenlines. The scalar $Q_0^*$ is nonzero on these factors.

In Schoen's notation, $W_0$ is the quotient of $Z^4$ by
the sum-zero deck subgroup together with $\mathfrak S_4$.
Its degree-three map to $\operatorname{Sym}^4C$ is etale.
The inverse image of $Q$ is the disjoint union of three
copies of $\mathbb P^2$. The nontrivial character sums
of their classes are nonzero by the normal-bundle
intersection in Schoen's Lemma 2.6. Pulling them back
to $Z^4$ gives nonzero classes in $U$ represented by
character sums of the ordered inverse images of $T_j$.
This is precisely the unramified case of Theorem 2.0,
and does not use an inverse Lefschetz correspondence.

Since $f_*f^*=(\deg f)\operatorname{Id}$, the pushforward
of either nonzero class in (F1S10.4) is a nonzero Weil class.
On a connected norm torsor the endomorphism $Q_0$ in
(F1S10.3) is multiplication by three followed by a
translation: its norm term is constant there.
The ordered quotient has generic degree $4!=24$, and
$T_j\to S_j$ has generic degree one by HD-F1-T07.
Consequently the pushed character sum is the corresponding
nonzero sum of $24[3]_*[S_j]$. Deck transformation permutes
the three surfaces and acts invertibly on $W_K$.
Thus every $[S_j]$ has a nonzero Weil component.

Equation (F1S10.2), the degree $\beta^2[S_j]=36$, and
$\beta^4=4!\cdot9=216$ now give
\[
                    [S_j]=\beta^2/6+w_j,\qquad 0\ne w_j\in W_K.
 \tag{F1S10.5}
\]
The product $\beta^2w_j$ vanishes, either by the determinant
description or by the unitary eigenbasis below.

\paragraph{2. Positivity determines the exact pencil count.}
Suppress the index $j$ and let $r$ be its number of
exceptional pencils. Adjunction for the smooth embedding
$S\subset P$ gives
\[
 [S]^2=\int_S c_2(N_{S/P})=K_S^2-c_2(S)=6+2r.
\]
Combining with (F1S10.5) yields
\[
                              w^2=2r.
 \tag{F1S10.6}
\]
Choose a $\beta$-unitary basis $a_1,a_2,b_1,b_2$ of
$H^{1,0}(P)$ adapted to the two $K$ eigenspaces, so
\[
 \beta=\sqrt{-1}\sum_{i=1}^4z_i\wedge\bar z_i,
 \qquad(z_1,z_2,z_3,z_4)=(a_1,a_2,b_1,b_2).
\]
Put $A=a_1\wedge a_2$, $B=b_1\wedge b_2$. Then
\[
 \int_P A\wedge B\wedge\bar A\wedge\bar B
        =\beta^4/24=9,\qquad
 w=\lambda A\wedge\bar B+\bar\lambda B\wedge\bar A.
 \tag{F1S10.7}
\]
In particular $w^2=18|\lambda|^2$.
The restriction pairing on the holomorphic two-forms
spanned by $A,B$ is exactly
\[
 \left(\int_P\Omega_i\wedge\overline{\Omega_j}\wedge[S]\right)
       _{\Omega_1=A,\Omega_2=B}
   =\begin{pmatrix}3&9\bar\lambda\\9\lambda&3\end{pmatrix}.
 \tag{F1S10.8}
\]
For example the diagonal entries are $2\cdot9/6=3$;
the off-diagonal entry comes from the complementary
determinant summand in (F1S10.7). These are actual
integrals of restricted forms on the smooth surface.
The matrix is positive semidefinite, because
$\int_S\omega\wedge\bar\omega\geq0$ for a holomorphic
two-form, with equality precisely when that form vanishes.
Thus $|\lambda|\leq1/3$ and $w^2\leq2$.
Since $w\ne0$ and (F1S10.6) has integral $r\geq0$,
\[
                    r=1,\qquad w^2=2,\qquad |\lambda|=1/3.
 \tag{F1S10.9}
\]
The resulting invariants are $K_S^2=16$ and $c_2(S)=8$.
The form
\[
                         \Omega=A-3\bar\lambda B
 \tag{F1S10.10}
\]
has zero restriction by (F1S10.8). Its square is nonzero,
so it is a nondegenerate invariant two-form on $P$.
Hence $S$ is Lagrangian for this specified form.

\paragraph{3. The irregularity is four.}
The embedding in the simple abelian fourfold gives
$q(S)\geq4$: a one-form vanishing on $S$ vanishes on
the abelian variety generated by $S$, which is $P$.
We show that $S$ has no fibration onto a curve of genus
at least two. Let $f:S\to B_0$ be such a fibration with
connected fibres, and write $b=g(B_0)\geq2$, $g=g(F)$
for the genus of its smooth fibre. The curve $F$ generates
$P$ by simplicity, so the induced map $J(F)\to P$ is
surjective and $g\geq4$.

Here is the Euler inequality with its hypotheses verified.
The cotangent bundle $\Omega_S^1$ is a quotient of the
trivial bundle $\Omega_P^1|_S$. Thus $K_S$ is globally
generated, in particular nef. For a fibre divisor
$D=\sum_i m_iC_i$, one has $D\cdot C_i=0$. Writing
$D_{\mathrm{red}}=\sum_iC_i$ gives
\[
 D_{\mathrm{red}}^2
   =-\sum_{i<j}(C_i\cdot C_j)
       \left(\sqrt{m_i/m_j}-\sqrt{m_j/m_i}\right)^2\leq0.
\]
For any connected reduced curve,
$e(D_{\mathrm{red}})\geq2-2p_a(D_{\mathrm{red}})$.
Indeed normalization gives this inequality from
$2\delta_x\geq r_x-1$ at every singular point with
$r_x$ branches. Adjunction and nefness now imply
\[
 -e(D_{\mathrm{red}})
 \leq K_SD_{\mathrm{red}}+D_{\mathrm{red}}^2
 \leq K_SD_{\mathrm{red}}\leq K_SD=2g-2.
\]
All singular-fibre Euler corrections are consequently
nonnegative. Additivity of Euler characteristic over
the smooth locus of the base and its finite complement
gives
\[
 8=c_2(S)=4(b-1)(g-1)+
             \sum_{s\in B_0}\bigl(e(D_{s,\mathrm{red}})+2g-2\bigr)
       \geq12,
\]
a contradiction.

Now $\chi(\mathcal O_S)=2$ gives $p_g=q+1$. Put
$m=\binom q2$. The kernel of
\[
 \bigwedge^2H^0(S,\Omega_S^1)\longrightarrow H^0(S,K_S)
 \tag{F1S10.11}
\]
has dimension at least $m-q-1$. For $q\geq5$ its
projectivization intersects $\operatorname{Gr}(2,q)$,
because the sum of their dimensions minus $m-1$ is
at least $q-5\geq0$. Hence two independent holomorphic
one-forms have zero wedge. The Castelnuovo--de Franchis
argument then supplies a fibration onto a curve of
genus at least two, already excluded.
Explicitly, the meromorphic ratio of the two forms is
nonconstant, and closedness gives $d(\omega_1/\omega_2)
\wedge\omega_2=0$. After resolving this map and taking
Stein factorization, both forms descend to the resulting
curve. They have no poles there: pullback of a pole of
a one-form under $t=u^m$ still has a pole. Their independence
forces genus at least two. Exceptional rational curves
map constantly to that curve, so the fibration descends
to $S$ itself. Thus $q\leq4$.
We obtain
\[
                         q=4,\quad p_g=5,\quad h^{1,1}(S)=12.
 \tag{F1S10.12}
\]
The last equality follows from $c_2=2-4q+2p_g+h^{1,1}$.

\paragraph{4. The exact conormal map has rank twelve.}
The unitary change of basis
$a_1'=e^{\sqrt{-1}\arg\lambda}a_1$,
$a_2'=a_2$, $b_i'=b_i$ preserves $\beta$ and the two
$K$ eigenspaces. In this basis (F1S10.7) reads
\[
                         w=(A\bar B+B\bar A)/3.
 \tag{F1S10.13}
\]
Thus this is an actual invertible change of coordinates
in the exterior algebra, preserving both the real structure
and the polarization. It identifies every norm-two case
here with the following single intersection calculation.

For a matrix $M\in\operatorname{Mat}_4(\mathbb C)$ put
$\alpha_M=\sum M_{ij}z_i\wedge\bar z_j$, and write it
in $2\times2$ blocks $M=\left(\begin{smallmatrix}A_0&B_0\\C_0&D_0
\end{smallmatrix}\right)$. For another matrix $M'$ use
primed blocks, and put
$d(U,V)=\operatorname{tr}(U)\operatorname{tr}(V)-\operatorname{tr}(UV)$,
the polarization of the $2\times2$ determinant.
Direct wedge multiplication using (F1S10.13) gives
\[
 \int_P\alpha_M\alpha_{M'}[S]
 =-3\left(\operatorname{tr}(M)\operatorname{tr}(M')
          -\operatorname{tr}(MM')+d(B_0,B_0')+d(C_0,C_0')\right).
 \tag{F1S10.14}
\]
This formula also determines its radical without a rank
assertion based on a numerical approximation. Varying
the diagonal blocks of $M'$ gives
$A_0=D_0=(\operatorname{tr}A_0+\operatorname{tr}D_0)I_2$,
which forces $A_0=D_0=0$. Varying the other blocks gives
$C_0=\operatorname{adj}(B_0)$ and
$B_0=\operatorname{adj}(C_0)$; the latter is redundant
for $2\times2$ matrices. The radical is therefore exactly
\[
 \left\{\begin{pmatrix}0&B_0\\\operatorname{adj}(B_0)&0\end{pmatrix}:
                         B_0\in\operatorname{Mat}_2(\mathbb C)\right\},
 \tag{F1S10.15}
\]
of dimension four. The Gram form has rank twelve.
The independent exact exterior computation
\texttt{followup\_1/check\_norm\_surface\_pairings.py}
and its JSON certificate verify this full $16\times16$
matrix, its four-dimensional radical, and the numerical
intersections of HD-F1-S09.

The Gram form factors as restriction to $H^{1,1}(S)$
followed by its Poincare pairing. By (F1S10.12) the latter
space has dimension twelve. Hence the restriction
\[
                         H^1(P,\Omega_P^1)\longrightarrow
                         H^1(S,\Omega_S^1)
 \tag{F1S10.16}
\]
is surjective.

We finish by checking that (F1S10.16) is the requested
normal map, with its actual bundle identifications.
Contraction by the nondegenerate form (F1S10.10) gives
$T_P\simeq\Omega_P^1$. Since $\Omega|_S=0$, it also
gives $N_{S/P}\simeq\Omega_S^1$, and the square
formed by these isomorphisms, the quotient
$T_P|_S\to N_{S/P}$, and restriction
$\Omega_P^1|_S\to\Omega_S^1$ commutes.
Taking $H^1$ identifies (F1S10.16) with
$H^1(P,T_P)\to H^1(S,N_{S/P})$.
HD-F1-S01 identifies this latter map with
\[
 H^1(P,\mathcal O_P)\otimes\bigwedge^3H^0(P,\Omega_P^1)
       \longrightarrow H^1(S,N_{S/P}^{\vee}\otimes K_S).
 \tag{F1S10.17}
\]
Its source has dimension sixteen and its target twelve.
This proves the rank, kernel dimension and zero cokernel
in (F1S10.1), for the specified original norm-image surface.

\paragraph{5. Bloch's map and the actual family detector.}
By Serre duality (F1S10.17) is dual to Bloch's
semiregularity map; hence that map is injective.
For a deformation class $\kappa$ of $P$, the compatibility
of the normal obstruction and the cycle-class map is
\[
                  \sigma_{\mathrm{Bloch}}(\operatorname{ob}_S\kappa)
                                =\kappa\mathbin{\lrcorner}[S].
 \tag{F1S10.18}
\]
The whole class (F1S10.5), with its flat determinant
component $w$, remains Hodge on the four-dimensional
Weil domain. Thus its tangent space is contained in
$\ker\operatorname{ob}_S$. The calculation above shows
that this kernel itself has dimension four, so these
two tangent spaces are equal.

Ran's small-extension semiregularity theorem, in the
precise form already checked in HD-F1-S05, puts every
relative embedded obstruction over the full-class Hodge
locus in the kernel of Bloch's map. That kernel is zero
here. The relative Hilbert morphism is therefore smooth
at the actual point $[S\subset P]$ and has open image
in this Weil family. Its universal subscheme supplies
actual nearby smooth embedded detectors of degree $36$.
This is the family conclusion of the computed map,
and it preserves the full class $\beta^2/6+w$.

\paragraph{6. Extension to every nonhyperelliptic cover.}
Work over a fine level cover of the moduli space of smooth
nonhyperelliptic genus-three curves and then over a connected
component $\mathcal H$ of its finite etale cover parametrizing
a nonzero three-torsion bundle. Every such component maps
surjectively to the relevant connected curve-moduli component:
a finite etale map has image both open and closed.
The rank-four Prym differential used in paragraph 1 exists
on every $\mathcal H$. Indeed the failure of base-point
freeness of either $K_C\otimes\epsilon^{\pm1}$ implies
$\epsilon=\mathcal O(q-p)$ for distinct points $p,q$.
The relation $\epsilon^3=\mathcal O$ gives a degree-three
function with divisor $3q-3p$. Its two totally ramified
fibres leave at most six other branch values. After
quotienting by the one-dimensional automorphism group
preserving the first two values, these covers have
moduli dimension at most five. The finite-torsion
space $\mathcal H$ has dimension six, so this bad locus
cannot contain a component. This is the exact dimension
argument in Schoen (1988), following Lemma 3.11.
Each $\mathcal H$ therefore has very general points to
which paragraphs 1--5 apply.

The relative symmetric power of the smooth covering
curve is smooth over $\mathcal H$. Its finite map to
the relative symmetric power of the base curve is flat
of degree $81$; this can also be checked fibrewise by
the same Cohen--Macaulay argument as in HD-F1-S09.
The inverse image of the relative canonical plane
bundle is consequently flat and proper over
$\mathcal H$. The differential proof of HD-F1-S09
applies on every fibre. It follows that the resulting
family $\mathcal T\to\mathcal H$ is smooth and proper.
It has three connected components on every fibre,
permuted cyclically by the deck transformation.
After a finite etale local change these components
are individually labelled.

The first Betti number of a smooth proper complex
family is locally constant. Hence
$q(T_0)+q(T_1)+q(T_2)$ is locally constant; it equals
$12$ at the very general fibres already calculated.
The deck isomorphisms make its three summands equal.
Thus every $T_j$ in this family has $q=4$. Blowdown
preserves irregularity, so every actual Abel image
$S_j$ has $q=4$ as well.

There is also a flat cohomology class throughout this
family, without any assumption about flatness of its
image subschemes. Namely the proper relative Abel
map pushes the fundamental class of a labelled
$T_j$ into the fourth cohomology of its connected
norm torsor. These pushforwards commute with
Gauss--Manin transport. One can check this directly
by pairing with any parallel degree-four test class:
the value is its integral after pullback to the
smooth proper family $T_j$, which is locally
constant. Translating a norm torsor to its Prym
does not change cohomology. Finally the generic
degree-one statement in HD-F1-T07 identifies this
pushforward with $[S_j]$ on every fibre.

The polarization class and the determinant subsystem
$W_K$ are parallel as well. Thus the already established
generic identity
\[
                [S_j]-\beta^2/6\in W_K,\qquad
                ([S_j]-\beta^2/6)^2=2
 \tag{F1S10.19}
\]
continues along every path in $\mathcal H$ (with a
component label transported along that path).
Membership in a flat subsystem and this cup-product
value are preserved by transport. No new projection
onto an enlarged special-fibre Hodge space is being
made here: the entire cohomology class is transported.

For an arbitrary fibre, HD-F1-S09 and HD-F1-T07 still
give a smooth embedded image and the identities
$K_S^2=15+r$, $c_2(S)=9-r$. Substitution of (F1S10.19)
into the self-intersection formula again gives
$2r=w^2=2$. Hence $r=1$, $K_S^2=16$, and $c_2(S)=8$
on that fibre. Together with the transported $q=4$
and $\chi=2$, this gives $h^{1,1}=12$ there.
The invariant two-form (F1S10.10), the unitary change
of coordinates, and the Gram calculation
(F1S10.14)--(F1S10.17) now apply verbatim to the
transported full class. They prove the surjectivity
of the requested conormal map on this arbitrary
nonhyperelliptic cover. Paragraph 5 then applies
at its actual embedded surface as well.

Every complex cover in the statement is represented
after a choice of fine level and a local family of
this form. This proves (F1S10.1) with the full scope
claimed at the start.

\paragraph{7. The principal degree-thirty-six detector from the original construction.}
There is now a route from the original norm image to
the principal detector in HD-F1-S07 that uses the
actual maps proved above. The ambient degeneration
of HD-F1-S08 identifies the polarized $K$ model at
its compact-type limit with $J\times J$, with
\[
 B_0=\begin{pmatrix}2&-1\\-1&2\end{pmatrix},\qquad
 T_0=\begin{pmatrix}-1&1\\-1&0\end{pmatrix}.
\]
Here $J$ is a principally polarized genus-two Jacobian.
Varying $J$ in the connected principally polarized
surface family reaches the product $E_0^2$, where
$E_0$ has Eisenstein multiplication. This locates
the original Prym family in the exact unitary
component used by HD-F1-S07.

With $\zeta^2+\zeta+1=0$, that section's actual
elliptic-block isogeny is
\[
 R=\begin{pmatrix}1&\zeta\\1&\zeta^2\end{pmatrix},\qquad
 R^\dagger R=B_0,\qquad
 RT_0=\operatorname{diag}(\zeta^2,\zeta)R.
\]
The map $R\oplus R$ has degree nine and gives a
$K$-stable superlattice on which the descended
polarization is principal. The same superlattice
over the common marked period domain gives the
polarized $K$-linear quotient
\[
                   q:P\longrightarrow A,\qquad q^*\theta=\beta,
                   \quad\deg q=9.
\]
After the finite level changes already described in
HD-F1-S07, it is an actual isogeny of families.
The map on the marked domains is the identity,
so its image is dense in the specified principal
Eisenstein component.

Take the actual norm-image surface $S\subset P$
proved above and let $Z=q(S)_{\mathrm{red}}$.
The generic degree of $q|_S$ is
$|\ker(q)\cap\operatorname{Stab}(S)|$: a translation
that identifies generic points of the integral
surface preserves the whole surface. This subgroup
has order dividing nine and acts freely on the
smooth $S$. Its quotient is finite etale, so its
order also divides $\chi(\mathcal O_S)=2$.
Consequently $q|_S$ has generic degree one. Thus
\[
 \deg_\theta Z=36,\qquad
 [Z]=\frac32\theta^2+w',\qquad
 w'=q_*w\in W_K(A),\qquad (w')^2=18.
 \tag{F1S10.20}
\]
The equalities use $q_*\beta^2=9\theta^2$ and
$q^*q_*w=9w$. Moreover $w'$ is integral:
$[Z]$ is integral and $\theta^2/2$ is integral
for a principal alternating form, as is seen
by expanding its symplectic-basis Pfaffians.
Properness of the relative Chow space carries
this entire flat class to an effective degree-36
cycle on every point of the principal component,
with the determinant component still of norm
eighteen. This proves the same uniform detector
input of HD-F1-S07 and HD-F1-T05 directly from
the original canonical norm construction.
